% Gerado por conferir_calculos.py — não editar à mão
\begin{longtable}{L{0.40\textwidth}L{0.20\textwidth}L{0.20\textwidth}c}
\caption{Casos especiais e propriedades verificadas pela implementação de referência}\label{tab:casos-especiais}\\
\hline
Caso & Esperado & Obtido & Situação \\
\hline
\endfirsthead
\hline
Caso & Esperado & Obtido & Situação \\
\hline
\endhead
SAC com i = 0: prestação constante = PV/n & 83,33 (todas) & 83,33 / juros 0,00 & OK \\
Price com i = 0: PMT = PV/n (limite) & 83,33; juros 0 & 83,33; juros 0,00 & OK \\
n = 1 (SAC e Price coincidem) & 1010,00 & 1010,00 e 1010,00 & OK \\
SAC: juros em PA de razão -i$\cdot$PV/n (C1) & -0,833333\ldots{} & -0,833333 & OK \\
SAC: total de juros = i$\cdot$PV$\cdot$(n+1)/2 (C1) & 6050,00 & 6050,00 & OK \\
Price: total de juros = n$\cdot$PMT $-$ PV (C1) & 7216,51 & 7216,51 & OK \\
Juros compostos: laço (Código~\ref{lst:mvc}) $\times$ fórmula, C=1000, i=1\%, n=12 & 1126,825030 & 1126.8250301319695 (E\_r 2.2e-16) & OK \\
Juros compostos com n = 0 & 1000,00 & 1000,00 & OK \\
Juros simples C=1000, i=1\%, n=12 & 1120,00 & 1120,00 & OK \\
Aportes: recorrência = fórmula fechada (M0=0, A=100, i=1\%, n=12) & 1268,250301 & 1268,250301 & OK \\
Aportes com i = 0: M\_n = M0 + nA & 1200,00 & 1200,00 & OK \\
Taxa equivalente: 1\% a.m. $\to$ a.a. & 12,6825\% & 12,6825\% & OK \\
Taxa equivalente: 12\% a.a. $\to$ a.m. & 0,9489\% & 0,9489\% & OK \\
Tempo para dobrar a 1\% a.m. & 70 meses & 70 & OK \\
Problema da Figura~\ref{fig:modulo4}: R\$ 1000 na gaveta, inflação 5\% & perda 47,62 (preços iniciais) / falta 50,00 (preços finais) & 47,62 / 50,00 & OK \\
Missão da Figura~\ref{fig:modulo_poupa}: mesada 150, gasto 20, meta 650 & 5 meses & 5 & OK \\
Poupança com meta não divisível (150, 20, 700) & 6 meses (arredondar para cima) & 6 & OK \\
Poupança com gasto = mesada & erro tratado (meta inatingível) & ValueError tratado & OK \\
Entrada ``1,5'' lida com \texttt{float()}, como na versão anterior do Código~\ref{lst:mvc} & ValueError $\to$ mensagem de erro ao aluno & ValueError & OK \\
Entrada ``1,5'' lida com \texttt{ler\_numero\_br()} & 1,5 & 1,5 & OK \\
Formatação de 1234,5 com o especificador \texttt{.2f} & R\$ 1.234,50 (padrão brasileiro) & 1234.50 (versão anterior) $\times$ R\$ 1.234,50 & OK \\
\hline
\end{longtable}
